% !TEX TS-program = xelatex
% !BIB program = bibtex
% !TeX spellcheck = ru_RU

% About magic macros see also
% https://tex.stackexchange.com/questions/78101/

% По умолчанию используется шрифт 14 размера.
% Если Вы не влезаете в лимит страниц и нужен 12-й шрифт,
% то уберите опцию [14pt]
\documentclass[14pt, russian]{matmex-diploma-custom}

% !TeX spellcheck = ru_RU
% !TEX root = vkr.tex
% Опциональные добавления используемых пакетов. Вполне может быть, что они вам не понадобятся, но в шаблоне приведены примеры их использования.
\usepackage{tikz} % Мощный пакет для создание рисунков, однако может очень сильно замедлять компиляцию
\usetikzlibrary{decorations.pathreplacing,calc,shapes,positioning,tikzmark}

% Библиотека для TikZ, которая генерирует отдельные файлы для каждого рисунка
% Позволяет ускорить компиляцию, однако имеет свои ограничения
% Например, ломает пример выделения кода в листинге из шаблона
% \usetikzlibrary{external}
% \tikzexternalize[prefix=figures/]

\newcounter{tmkcount}

\tikzset{
	use tikzmark/.style={
			remember picture,
			overlay,
			execute at end picture={
					\stepcounter{tmkcount}
				},
		},
	tikzmark suffix={-\thetmkcount}
}

\usepackage{booktabs} % Пакет для верстки "более книжных" таблиц, вполне годится для оформления результатов
% В шаблоне есть команда \multirowcell, которой нужен этот пакет.
\usepackage{multirow}
\usepackage{siunitx} % для таблиц с единицами измерений

\newcommand{\cd}[1]{\texttt{#1}}
\newcommand{\inbr}[1]{\left<#1\right>}

% Для названий стоит использовать \textsc{}
\newcommand{\OCaml}{\textsc{OCaml}}
\newcommand{\Cpp}{\textsc{C++}}
\newcommand{\C}{\textsc{C}}
\newcommand{\miniKanren}{\textsc{miniKanren}}
\newcommand{\BibTeX}{\textsc{BibTeX}}
\newcommand{\vsharp}{\textsc{V$\sharp$}}
\newcommand{\fsharp}{\textsc{F$\sharp$}}
\newcommand{\csharp}{\textsc{C$\sharp$}}
\newcommand{\GitHub}{\textsc{GitHub}}
\newcommand{\SMT}{\textsc{SMT}}

\newcolumntype{L}[1]{>{\raggedright\let\newline\\\arraybackslash\hspace{0pt}}m{#1}}
%\newcolumntype{C}[1]{>{\centering\let\newline\\\arraybackslash\hspace{0pt}}m{#1}}
\newcolumntype{R}[1]{>{\raggedleft\let\newline\\\arraybackslash\hspace{0pt}}m{#1}}

%  Команды и пакеты, не используемые в шаблоне, которые тем не менее могут быть полезными.

% \newcolumntype{Y}{>{\centering\arraybackslash}X}

% \usepackage{mathrsfs}

% \lstdefinelanguage{ocaml}{
% keywords={@type, function, fun, let, in, match, with, when, class, type,
% nonrec, object, method, of, rec, repeat, until, while, not, do, done, as, val, inherit, and,
% new, module, sig, deriving, datatype, struct, if, then, else, open, private, virtual, include, success, failure,
% lazy, assert, true, false, end},
% sensitive=true,
% commentstyle=\small\itshape\ttfamily,
% keywordstyle=\ttfamily\bfseries, %\underbar,
% identifierstyle=\ttfamily,
% basewidth={0.5em,0.5em},
% columns=fixed,
% fontadjust=true,
% literate={->}{{$\to$}}3 {===}{{$\equiv$}}1 {=/=}{{$\not\equiv$}}1 {|>}{{$\triangleright$}}3 {\\/}{{$\vee$}}2 {/\\}{{$\wedge$}}2 {>=}{{$\ge$}}1 {<=}{{$\le$}} 1,
% morecomment=[s]{(*}{*)}
% }

\usepackage{totcount}
\usepackage{setspace}


\begin{document}
% TODO: Formatting
% !TeX spellcheck = ru_RU
% !TEX root = vkr.tex

%% Если что-то забыли, при компиляции будут ошибки Undefined control sequence \my@title@<что забыли>@ru
%% Если англоязычная титульная страница не нужна, то ее можно просто удалить.
\filltitle{ru}{
    %% Актуально только для курсовых/практик. ВКР защищаются не на кафедре а в ГЭК по направлению,
    %%   и к моменту защиты вы будете уже не в группе.
    chair              = {Кафедра, на которой работает научник},
    group              = {ХХ.БХХ-мм},
    %
    %% Макрос filltitle ненавидит пустые строки, поэтому обязателен хотя бы символ комментария на строке
    %% Актуально всем.
    title              = {Шаблон отчёта по учебной практике},
    %
    %% Здесь указывается тип работы. Возможные значения:
    %%   production - производственная практика;
    %%   coursework - отчёт по курсовой работе;
    %%   practice - отчёт по учебной практике;
    %%   prediploma - отчёт по преддипломной практике;
    %%   master - ВКР магистра;
    %%   bachelor - ВКР бакалавра.
    type               = {practice},
    %
    %% Здесь указывается вид работы. От вида работы зависят критерии оценивания.
    %%   solution - <<Решение>>. Обучающемуся поручили найти способ решения проблемы в области разработки программного обеспечения или теоретической информатики с учётом набора ограничений.
    %%   experiment - <<Эксперимент>>. Обучающемуся поручили изучить возможности, достоинства и недостатки новой технологии, платформы, языка и т. д. на примере какой-то задачи.
    %%   production - <<Производственное задание>>. Автору поручили реализовать потенциально полезное программное обеспечение.
    %%   comparison - <<Сравнение>>. Обучающемуся поручили сравнить несколько существующих продуктов и/или подходов.
    %%   theoretical - <<Теоретическое исследование>>. Автору поручили доказать какое-то утверждение, исследовать свойства алгоритма и т.п., при этом не требуя написания кода.
    kind               = {solution},
    %
    author             = {ФАМИЛИЯ Имя Отчество},
    %
    %% Актуально только для ВКР. Указывается код и название направления подготовки. Типичные примеры:
    %%   02.03.03 <<Математическое обеспечение и администрирование информационных систем>>
    %%   02.04.03 <<Математическое обеспечение и администрирование информационных систем>>
    %%   09.03.04 <<Программная инженерия>>
    %%   09.04.04 <<Программная инженерия>>
    %% Те, что с 03 в середине --- бакалавриат, с 04 --- магистратура.
    specialty          = {02.03.03 <<Математическое обеспечение и администрирование информационных систем>>},
    %
    %% Актуально только для ВКР. Указывается шифр и название образовательной программы. Типичные примеры:
    %%   СВ.5006.2017 <<Математическое обеспечение и администрирование информационных систем>>
    %%   СВ.5162.2020 <<Технологии программирования>>
    %%   СВ.5080.2017 <<Программная инженерия>>
    %%   ВМ.5665.2019 <<Математическое обеспечение и администрирование информационных систем>>
    %%   ВМ.5666.2019 <<Программная инженерия>>
    %% Шифр и название программы можно посмотреть в учебном плане, по которому вы учитесь.
    %% СВ.* --- бакалавриат, ВМ.* --- магистратура. В конце --- год поступления (не обязательно ваш, если вы были в академе/вылетали).
    programme          = {СВ.5006.2019 <<Математическое обеспечение и администрирование информационных систем>>},
    %
    %% Актуально только для ВКР, только для матобеса и только 2017-2018 годов поступления. Указывается профиль подготовки, на котором вы учитесь.
    %% Названия профилей можно найти в учебном плане в списке дисциплин по выбору. На каком именно вы, вам должны были сказать после второго курса (можно уточнить в студотделе).
    %% Вот возможные вариканты:
    %%   Математические основы информатики
    %%   Информационные системы и базы данных
    %%   Параллельное программирование
    %%   Системное программирование
    %%   Технология программирования
    %%   Администрирование информационных систем
    %%   Реинжиниринг программного обеспечения
    % profile            = {Системное программирование},
    %
    %% Актуально всем.
    %% Должно умещаться в одну строчку, допускается использование сокращений, но без переусердствования,
    %% короткая строка с большим количеством сокращений выглядит странно
    %supervisorPosition = {проф. кафeдры системного программирования, д.ф.-м.н.,}, % Терехов А. Н.
    %supervisorPosition = {ст. преподаватель кафедры ИАС, к.~ф.-м.~н. (если есть),}, % Смирнов К. К.
    supervisorPosition = {доцент кафедры системного программирования, к.~ф.-м.~н. (если есть),},
    supervisor         = {Научник~Н.~Н.},
    %
    %% Актуально только для практик и курсовых. Если консультанта нет, закомментировать или удалить вовсе.
    consultantPosition = {должность, ООО <<Место работы>>, степень  (если есть),},
    consultant         = {Консультант~К.~К.},
    %
    %% Актуально только для ВКР.
    reviewerPosition   = {должность, ООО <<Место работы>>, степень (если есть),},
    reviewer           = {Рецензент~Р.~Р.},
}

% \filltitle{en}{
%     chair              = {Advisor's chair},
%     group              = {ХХ.BХХ-mm},
%     title              = {Template for SPbU qualification works},
%     type               = {practice},
%     author             = {FirstName Surname},
%     %
%     %% Possible choices:
%     %%   02.03.03 <<Software and Administration of Information Systems>>
%     %%   02.04.03 <<Software and Administration of Information Systems>>
%     %%   09.03.04 <<Software Engineering>>
%     %%   09.04.04 <<Software Engineering>>
%     %% Те, что с 03 в середине --- бакалавриат, с 04 --- магистратура.
%     specialty          = {02.03.03 ``Software and Administration of Information Systems''},
%     %
%     %% Possible choices:
%     %%   СВ.5006.2017 <<Software and Administration of Information Systems>>
%     %%   СВ.5162.2020 <<Programming Technologies>>
%     %%   СВ.5080.2017 <<Software Engineering>>
%     %%   ВМ.5665.2019 <<Software and Administration of Information Systems>>
%     %%   ВМ.5666.2019 <<Software Engineering>>
%     programme          = {СВ.5006.2019 ``Software and Administration of Information Systems''},
%     %
%     %% Possible choices:
%     %%   Mathematical Foundations of Informatics
%     %%   Information Systems and Databases
%     %%   Parallel Programming
%     %%   System Programming
%     %%   Programming Technology
%     %%   Information Systems Administration
%     %%   Software Reengineering
%     % profile            = {Software Engineering},
%     %
%     %% Note that common title translations are:
%     %%   кандидат наук --- C.Sc. (NOT Ph.D.)
%     %%   доктор ... наук --- Sc.D.
%     %%   доцент --- docent (NOT assistant/associate prof.)
%     %%   профессор --- prof.
%     supervisorPosition = {Sc.D, prof.},
%     supervisor         = {S.S. Supervisor},
%     %
%     consultantPosition = {position at ``Company'', degree if present},
%     consultant         = {C.C. Consultant},
%     %
%     reviewerPosition   = {position at ``Company'', degree if present},
%     reviewer           = {R.R. Reviewer},
% }

\maketitle
\setcounter{tocdepth}{2}
\tableofcontents

% В курсаче/практике/ВКР не нужен
% \begin{abstract}
% \end{abstract}

\pagebreak
\newfontfamily\myfont{CMU Sans Serif}
\begin{center}
  \hspace{0pt}
  \vfill
  {\Huge\myfont
    Текст ВКР или учебной практики пишется не ради зачета, а чтобы люди его прочитали, поняли как круто Вы все сделали, и могли продолжить с того места, где Вы остановились.
  \vspace{2em}

  Повторять эту страницу в тексте вашей работы \emph{нельзя}.}
  \vfill
  \hspace{0pt}
\end{center}
\pagebreak

% !TeX spellcheck = ru_RU
% !TEX root = vkr.tex

\section*{Введение}
\thispagestyle{withCompileDate}

В современном мире существует множество инструментов, реализованных на разных языках программирования. Для того чтобы при разработке чего-либо на каком-либо языке программирования использовать надёжные готовые решения, написанные на других языках, реализуют привязки (от англ. bindings) --- интерфейс, предоставляющий набор функций для взаимодействия с библиотеками, написанными на других языках программирования.

LLVM \cite{LLVM} --- инфраструктура для разработки компиляторов. LLVM состоит из множества подпроектов. Ядро LLVM базируется на промежуточном представлении LLVM (LLVM IR). LLVM IR это довольно низкоуровневое промежуточное представление, поэтому прямолинейное моделирование на нём высокоуровневых языковых конструкций затруднительно. Чтобы решить эту проблему, а также внедрить специфичные для языка оптимизации добавляют множество уровней абстракции, определяя новые промежуточные представления.

Разработка промежуточных представлений -- хорошо изученная область. Однако стоимость проектирования и внедрения промежуточных представлений всё равно остаётся высокой.

Проект MLIR \cite{MLIR} (Multi-Level Intermediate Representation) --- это подпроект LLVM, направленный на непосредственное решение проблем проектирования и реализации промежуточных представлений за счёт удешевления определения и внедрения новых уровней абстракции и предоставления инфраструктуры для решения распространенных проблем разработки промежуточных представлений.

Проблема в том, что MLIR реализован на языке \Cpp{}, но разработчикам, использующим \OCaml{}, также хотелось бы иметь такой мощный инструмент у себя в арсенале. Таким образом, эта работа посвящена реализации привязок к MLIR для языка программирования \OCaml{}.

% !TeX spellcheck = ru_RU
% !TEX root = vkr.tex

\section{Постановка задачи}
\label{sec:task}
Целью работы является реализация привязок к MLIR для \OCaml{}. Для её выполнения были поставлены следующие задачи:

\begin{enumerate}
	\item реализовать привязки к MLIR;
	\item разработать компилятор для мини языка программирования с использованием разработанных привязок;
	\item реализовать тесты.
\end{enumerate}

% !TeX spellcheck = ru_RU
% !TEX root = vkr.tex

\section{Обзор}
\label{sec:relatedworks}

\subsection{SSA форма}

Промежуточное представление соответствует SSA \cite{SSA} форме тогда и только тогда, когда каждой переменной значение присваивается ровно один раз.

SSA форма используется во многих современных компиляторах императивных языков программирования, например, в LLVM, HotSpot JVM и GCC. Использование SSA формы значительно упрощает анализ потока управления, поэтому внедрение многих оптимизаций (например, свёртка констант, удаление неиспользуемого кода и удаление общих подвыражений) становится легче.

\subsection{MLIR}

MLIR \cite{MLIR} (Multi-Level Intermediate Representation) --- это гибкая инфраструктура для разработки промежуточных представлений компиляторов на C++. MLIR имеет встроенное (от англ. builtin) промежуточное представление, которое соответствует SSA форме и содержит базовые операции, типы и атрибуты. Из базового промежуточного представления можно создавать собственные промежуточные представления, также в SSA форме, добавляя к базовому промежуточному представлению новые типы и операции, которые могут иметь любую семантику, заданную пользователем.

Далее обзор базовых возможностей MLIR.

\subsubsection{Базовые объекты MLIR}

\begin{itemize}
	\item
	      \textbf{Операции}

	      Операция --- это основная сущность в MLIR. Всё от инструкций до модулей и функций представляется как операция. MLIR позволяет пользователям определять операции.

	      Операция имеет уникальное имя, состоящее из названия диалекта, которому она принадлежит, и идентификатора. Операция имеет нуль или более операндов и результатов. Кроме того, операция может содержать атрибуты, регионы и информацию о местоположении.

	      \begin{lstlisting}[caption={Пример операции.}, frame=single, label={op}]
%res:2 = "mydialect.op"(%arg1, %arg2) { attribute1 = true,
    attribute2 = 1.5} : (tensor<*xf64>, i32) -> (f32, i8)
                     loc(callsite("foo" at "mysrc.e":10:8))
		  \end{lstlisting}

	      На листинге \ref{op} определена абстрактная операция <<op>> из диалекта <<mydialect>>, которая имеет два операнда <<\%arg1>> и <<\%arg2>> типа \texttt{tensor$<$*xf64$>$} (безранговый тензор с элементами типа \texttt{float64}) и \texttt{int32} соответственно; два именованных атрибута: <<attribute1>> со значением \texttt{true}, <<attribute2>> со значением \texttt{1.5}; два результата типа \texttt{float32} и \texttt{int8} соответственно. А также операция имеет информацию о местоположении в исходном тексте программы.

	\item
	      \textbf{Атрибуты}

	      Атрибут --- это любая compile-time константа, которая прикреплена к операции. Любая операция имеет словарь именованных атрибутов, у которого ключ --- имя атрибута, значение --- атрибут. Пользователь может определять собственные атрибуты.

	\item
	      \textbf{Типы}

	      Каждое значение в MLIR имеет тип. Система типов в MLIR полностью расширяема.

	      MLIR предоставляет базовые типы. Например, целочисленные типы, типы с плавающей точкой, строки, кортежи, многомерные вектора, тензоры.

	\item
	      \textbf{Блоки и регионы}

	      Блок --- это список операций, который заканчивается операцией-терминатором. Блоки образуют граф потока управления программы (от англ. control flow graph). В MLIR для поддержания SSA формы вместо использования $\phi$-узлов (как например в LLVM IR) используются аргументы блоков (терминатор блока непосредственно передаёт значения в аргументы следующего блока).

	      Регион --- это список блоков. Любая операция может содержать регионы. При помощи регионов в MLIR реализована вложенная структура: регион содержит список блоков, а блоки содержат список операций, которые, в свою очередь, могут содержать регионы.

	\item
	      \textbf{Диалекты}

	      Диалект в MLIR --- это пространство имён (от англ. namespace) для операций, типов и атрибутов, при помощи которого MLIR поддерживает расширяемость.

	      Базовое промежуточное представление MLIR это не что иное, как диалект, который называется <<builtin>>\footnote{\url{https://mlir.llvm.org/docs/Dialects/Builtin/} (дата доступа:   \DTMdate{2024-01-04}).}.

	      В репозитории MLIR также содержится множество других диалектов\footnote{\url{https://mlir.llvm.org/docs/Dialects/} (дата доступа:   \DTMdate{2024-01-04}).}.
\end{itemize}

\subsubsection{TableGen}

MLIR стремится сделать определение новых диалектов и операций, а также правил переписывания (от англ. rewrite rules\footnote{\url{https://mlir.llvm.org/docs/DeclarativeRewrites/} (дата доступа:   \DTMdate{2024-01-05}).}) паттернов в IR (Intermediate Representation) простым. Для этого MLIR предоставляет декларативный DSL (Domain Specific Language) --- TableGen\footnote{\url{https://mlir.llvm.org/docs/DefiningDialects/Operations/} (дата доступа:   \DTMdate{2024-01-04}).}, из которого генерируется \Cpp{}.

Для того, чтобы определить свой диалект необходимо лишь определить класс, который наследуется от базового класса <<Dialect>>. Внутри класса определить имя, краткое описание, полное описание диалекта и имя для пространства имён диалекта в \Cpp{}.

\begin{lstlisting}[caption={Пример определения диалекта <<MyDialect>>, при помощи TableGen.}, frame=single]
def MyDialect : Dialect {
    let name = "myDialect";

    let summary = "My dialect.";
    let description = [{ My cool dialect. }];

    let cppNamespace = "myDialect";
}
\end{lstlisting}

\begin{lstlisting}[caption={Пример определения операции сложения тензоров, при помощи TableGen.}, label={op_tblgen}, frame=single, mathescape=true]
def TransposeOp : Op<MyDialect, "transpose", [Pure]> {
    let summary = "transpose operation";
    let description = [{
            The op takes tensor and returns same tensor
            with reversed shape.
    }];

    let arguments = (ins F64Tensor:$\mbox{\$}$arg);
    let results = (outs F64Tensor);

    let assemblyFormat = [{
        "(" $\mbox{\$}$arg ":" type($\mbox{\$}$arg) ")" attr-dict "to" type(results)
    }];
}
\end{lstlisting}

Для определения операции необходимо указать какому диалекту принадлежит операция, задать уникальное имя (в рамках диалекта), список признаков (от англ. traits), которые задают поведение операции, аргументы, атрибуты и результаты операции. Также операция может содержать краткое и полное описание, из которого можно сгенерировать документацию для операции (в формате Markdown). Кроме того, для операции можно определять красивый вывод (от англ. pretty print), указывать на наличие верификаторов операции, которые можно определять в \Cpp{}.

На листинге \ref{op_tblgen} приведён пример определения операции диалекта <<MyDialect>>, которая моделирует транспонирование тензоров с элементами типа \texttt{float64}. Операция называется <<transpose>> и имеет признак <<pure>>, означающий, что операция детерминирована и не имеет побочных эффектов. Кроме этого, операция имеет красивый вывод.

\begin{lstlisting}[caption={Пример вывода операции, определённой на листинге \ref{op_tblgen}, без переопределённого вывода.}, frame=single]
%1 = "myDialect.transpose"(%0):(tensor<2x3xf64>)->tensor<3x2xf64>
\end{lstlisting}

\begin{lstlisting}[caption={Пример вывода операции, определённой на листинге \ref{op_tblgen}, с переопределённым выводом.}, frame=single]
%1 = myDialect.transpose(%0 : tensor<2x3xf64>) to tensor<3x2xf64>
\end{lstlisting}

\textbf{Правила переписывания}

TableGen позволяет декларативно определять правила переписывания\footnote{\url{https://mlir.llvm.org/docs/DeclarativeRewrites/} (дата доступа:   \DTMdate{2024-01-05}).} шаблонов в IR.

\begin{lstlisting}[caption={Базовый класс для шаблона.}, frame=single]
class Pattern<
    dag sourcePattern, list<dag> resultPatterns,
    list<dag> additionalConstraints = [],
    list<dag> supplementalPatterns = [],
    dag benefitsAdded = (addBenefit 0)>;
\end{lstlisting}

Правило переписывания содержит две главные компоненты: исходный паттерн и итоговый паттерн. При применении правила исходный паттерн в IR будет заменён на итоговый паттерн.

\textbf{Пример}

Так как применение чётного количество раз операции транспонирования к тензору это тождественное преобразование, то можно определить правило переписывания, которое будет заменять вызов чётного числа транспонирования некоторого тензора на данный тензор.

\begin{lstlisting}[caption={Пример создания правила переписывания для операции, определённой на листинге \ref{op_tblgen}}, label={rewrite_pat}, frame=single]
// transpose(transpose(x)) -> x
def RedundantTransposeOptPattern :
  Pattern<(TransposeOp(TransposeOp $\mbox{\$}$arg)),
          [(replaceWithValue $\mbox{\$}$arg)]>;
\end{lstlisting}

Правило переписывания на листинге \ref{rewrite_pat} находит в IR двойной вызов операции <<myDialect.transpose>> от некоторого аргумента, захватывает аргумент первого вызова и заменяет эту последовательность на аргумент первого вызова.

\subsection{Взаимодействие OCaml с C}

Далее будет кратко рассказано как код, написанный на языке C можно вызывать из OCaml. Более подробно написано здесь\footnote{\url{https://v2.ocaml.org/manual/intfc.html} (дата доступа:   \DTMdate{2024-01-04}).}.

\subsubsection{Объявление}

Для того чтобы объявить внешнюю функцию, в \OCaml{} существует ключевое слово <<external>>.
\begin{lstlisting}[caption={Синтаксис объявления внешней функции.},language=Caml, frame=single]
external name : type = C-function-name
\end{lstlisting}

\begin{lstlisting}[caption={Пример объявления внешней функции <<print\_int>>, которая ссылается на функцию, определённую в С, с именем <<caml\_print\_int>>.},language=Caml, frame=single]
external print_int : int -> unit = "caml_print_int"
\end{lstlisting}

\subsubsection{Определение}

С функция должна принимать столько же аргументов, сколько и в объявлении внешней функции (если количество аргументов больше 5, то см. здесь\footnote{\url{https://v2.ocaml.org/manual/intfc.html\#ss:c-prim-impl} (дата доступа:   \DTMdate{2024-01-04}).}), типа \texttt{value} и возвращать значение типа \texttt{value}. Тип \texttt{value} служит для представления \OCaml{} типа.

\begin{lstlisting}[caption={Содержимое файла C с определением функции <<caml\_print\_int>>.}, language=C, frame=single]
#include <stdio.h>
#include <caml/mlvalues.h>

CAMLprim value caml_print_int(value val)
{
    int num = Int_val(val);
    printf("%d\n", num);
    return Val_unit;
}
\end{lstlisting}

Заголовочный файл <<caml/mlvalues.h>> предоставляет макросы для преобразования значений \OCaml{}-типа в тип \C{} и наоборот. Функция <<print\_int>>, определённая на стороне \OCaml{}, принимает \texttt{int} и возвращает \texttt{unit}, поэтому аргумент функции <<caml\_\-pri\-nt\-\_int>> типа \texttt{value} представляет из себя целое число. <<Int\_val>> преобразует значение типа \texttt{value} в значение С-типа \texttt{int}. Ожидается, что функция вернёт \texttt{unit}, поэтому функция <<caml\_\-pri\-nt\-\_int>> возвращает значение <<Val\_unit>>.

Таким образом, когда внешняя функция вызывается в \OCaml{}, вызывается функция \C{} с переданными аргументами. Затем значение, возвращаемое функцией, передается обратно в \OCaml{} как результат вызова функции.

Реализация внешней функции представляет из себя две задачи:
\begin{enumerate}
	\item декодирование аргументов для извлечения значений \C{} из значений \OCaml{} и кодирование возвращаемого значения как \OCaml{}-значения;
	\item фактическое вычисление какого-то результата из аргументов.
\end{enumerate}

Предпочтительно иметь две отдельные функции \C{} для реализации этих двух задач:

\begin{itemize}
	\item первая функция реализует основную логику функции, принимая и возвращая значения \C{};
	\item вторая --- часто называемая <<заглушкой>> (от англ. stub), представляет собой простую обёртку вокруг первой функции, которая преобразует свои аргументы из значений \OCaml{} в значения \C{}, вызывает первую функцию и преобразует возвращаемое значение \C{} в значение \OCaml{}.
\end{itemize}

\begin{lstlisting}[caption={Функция <<print\_int>> --- реализует логику, функция <<caml\_print\_int>> --- заглушка}, language=C, frame=single, label={qwe}]
void print_int(int num) {
   printf("%d\n", num);
   return;
}

CAMLprim value caml_print_int(value val)
{
    print_int(Int_val(num));
    return Val_unit;
}
\end{lstlisting}

\subsubsection{Линковка}

Существует два вида линковки \OCaml{} с \C{}:

\begin{enumerate}
	\item динамическая;
	\item статическая.
\end{enumerate}

Особенности динамической линковки:
\begin{enumerate}
	\item сохраняет независимость исполняемых файлов, скомпилированных в байт-код, от платформы;
	\item размер исполняемого файла меньше чем при статической линковки;
	\item пользователям библиотеки не нужно иметь компилятор и линковщик \C{}, а также runtime библиотеки \C{};
	\item итоговый исполняемый файл не является автономным (от англ. standalone): на машине должны быть библиотеки;
\end{enumerate}

\subsection{Ctypes}
\label{sec:ctypes}

Ctypes\footnote{\url{https://github.com/yallop/ocaml-ctypes} (дата доступа:   \DTMdate{2024-01-04}).} --- это \OCaml{} библиотека для создания привязок к \C{}. Основная цель библиотеки --- сделать создание привязок к \C{} максимально простым. Используя эту библиотеку, больше не нужно писать <<заглушки>> на С{}.

Так, например, для того чтобы вызвать из \OCaml{} ранее определённую на листинге \ref{qwe} функцию <<print\_int>>, необходимо предоставить имя функции, описать сигнатуру функции при помощи предоставляемых <<ctypes>> комбинаторов и слинковать \C{} и \OCaml{}.

\begin{lstlisting}[caption={Привязка к функции <<print\_int>>, определённой на листинге \ref{qwe}, при помощи <<ctypes>>}, language=Caml, frame=single, label ={ctypes_example}]
let print_int = foreign "print_int" (int @-> returning void)
\end{lstlisting}

\subsubsection{Линковка}

Существует два способа линковки привязок:

\begin{enumerate}
	\item динамическая;
	\item статическая генерация привязок.
\end{enumerate}

При динамической линковке ctypes использует библиотеку <<libffi>>\footnote{\url{https://github.com/libffi/libffi} (дата доступа:   \DTMdate{2024-01-04}).} для динамического открытия библиотек \C{}, поиска соответствующих символов вызываемой функции и упорядочивания аргументов функции в соответствии с ABI (Application Binary Interface) операционной системы. Многое из этого происходит <<под капотом>> ctypes, однако такой подход имеет следующие минусы:

\begin{itemize}
	\item понижение безопасности: библиотеки \C{} не содержат информацию о типах функций, которые они содержат, поэтому не существует способа проверить соответствуют ли переданные типы в <<foreign>> фактическим типам связываемой функции;
	\item накладные расходы на интерпретацию.
\end{itemize}

К счастью, ctypes предоставляет альтернативный способ линковки привязок --- статический. В этом случае линковка выполняется один раз во время сборки, вместо того, чтобы выполнять её при каждом вызове функции. При статической линковке:
\begin{enumerate}
	\item Используя те же описания функций, как на листинге \ref{ctypes_example}, генерируется .ml файл, который содержит объявления внешних функций, и .c файл с заглушками;
	\item .ml и .c компилируются и линкуются в одну библиотеку;
\end{enumerate}

\subsection{Существующие решения}

На просторах \GitHub{} была найдена попытка\footnote{\url{https://github.com/tachukao/ocaml-mlir} (дата доступа:   \DTMdate{2024-01-04}).} создания привязок к MLIR для \OCaml{}. Данные привязки создавались к форку\footnote{\url{https://github.com/tachukao/llvm-project} (дата доступа:   \DTMdate{2024-01-05}).} MLIR, базирующемуся на древней версии MLIR --- одиннадцатой, не примитивного примера использования привязок в репозитории не было, и на момент начала работы репозиторий был заброшен. Показалось, что использование некоторой базы может сильно сэкономить время, поэтому эти привязки были взяты за основу.

% !TeX spellcheck = ru_RU
% !TEX root = vkr.tex

\section{Background (опционально)}
Здесь пишется некоторая дополнительная информация о том, зачем делается то, что делается.

Например, в работе придумывается какой-то новый метод решения формул в \SMT{} в теориях с числами. Без каких-то дополнительных пояснений будет казаться, что работа состоит из жестокого \enquote{матана} и совсем не по теме кафедры системного программирования.
Поэтому, в данном разделе стоит рассказать, что все эти методы примеряются для верификации в проекте \vsharp{}, и поэтому непосредственно связаны с тематикой кафедры.

% !TeX spellcheck = ru_RU
% !TEX root = vkr.tex

\section{Метод}
Реализация в широком смысле: что таки было сделано. Скорее всего самый большой раздел.

\emph{Крайне желательно} к Новому году иметь что-то, что сюда можно написать.

Для понимания того как курсовая записка (отзыв по учебной практи\-ке/ВКР) должна писаться, можно посмотреть видео ниже. Они про научные доклады и написание научных статей, учебные практики и ВКР отличаются тем, что тут есть требования отдельных глав (слайдов) со списком задач и списком результатов. Но даже если вы забьёте на требования специфичные для ВКР, и соблюдете все рекомендации ниже, получившиеся скорее всего будет лучше чем первая попытка 99\% ваших однокурсников.

\begin{enumerate}
\item \enquote{Как сделать великолепный научный доклад} от Саймона Пейтона Джонса~\cite{SPJGreatTalk} (на английском).
\item \enquote{Как написать великолепную научную статью} от Саймона Пейтона Джонса~\cite{SPJGreatPaper} (на английском).
\item \enquote{Как писать статьи так, чтобы люди их смогли прочитать} от Дэрэка Драйера~\cite{DreyerYoutube2020} (на английском). По предыдующей ссылке разбираются, что должно быть в статье, т.е. как она должна выгля\-деть на высоком уровне. Тут более низкоуровнево изложено как должны быть устроены параграфы и т.п.
\item Ещё можно посмотреть How to Design Talks~\cite{JhalaYoutube2020} от Ranjit Jhala, но я думаю, что первых трех ссылок всем хватит.
\end{enumerate}

\subsection{Первая задача}
\label{subsec:task1}

\subsection{Вторая задача}
\label{subsec:task2}

\subsection{Третья задача}
\label{subsec:task3}


\subsection{Некоторые типичные ошибки}
Здесь мы будем собирать основные ошибки, которые случаются при написании текстов.
В интернетах тоже можно найти коллекции типич\-ных косяков\footnote{\href{https://www.read.seas.harvard.edu/~kohler/latex.html}{https://www.read.seas.harvard.edu/\textasciitilde kohler/latex.html} (дата доступа:   \DTMdate{2022-12-16}).}.

Рекомендуется по-умол\-ча\-нию использовать красивые греческие бук\-вы $\sigma$  и $\phi$, а именно $\phi$ вместо $\varphi$. В данном шаблоне команды для этих букв переставлены местами по сравнению с ванильным \TeX'ом.

Также, если работа пишется на русском языке, необходимо, чтобы работа была написана на \textit{грамотном} русском языке даже, если автор, из ближнего зарубежья\footnote{
Теоретически, возможен вариант написания текстов на английском языке, но это необходимо обсудить в первую очередь с научным руководителем.}.
Это включает в себя:
\begin{itemize}
  \item разделы должны оформляться с помощью \verb=\section{...}=, а также \verb=\subsection= и т.~п.; не нужно пытаться нумеровать вручную;
  \item точки после окончания предложений должны присутствовать;
  \item пробелы после запятых  и точек, в конце слова перед скобкой;
  \item неразрывные пробелы, там, где нужны пробелы, но переносить на другую строку нельзя, например \verb=т.~е.= или \verb=А.~Н.~Терехов=;
  \item дефис, там где нужен дефис (обозначается с помощью одиночного \enquote{минуса} (англ. dash) на клавиатуре);
  \item двойной дефис, там где он нужен; а именно  при указании проме\-жутка в цифрах: в 1900--1910 г.~г., стр. 150--154;
  \item даты стоит писать везде одинаково; чтобы об этом не следить, можно пользоваться заклинанием \verb=\DTMdate{2022-12-16}=;
  \item тире (т.~е. \verb=---= --- тройной минус) на месте тире, а не что-то другое; в русском языке тире не может \enquote{съезжать} на новую строку, поэтому стоит использовать такой синтаксис: \verb=До~--- после=;
  \item правильные двойные кавычки должны набираться с помощью пакета \texttt{csquotes}: для основного языка в \texttt{polyglossia} стоит использовать команду \verb=\enquote{текст}=, для второго языка стоит использовать \verb=\foreignquote{язык}{текст}=;
  \item все перечисления должны оформляться с помощью \verb=\enumerate= или \verb=\itemize=; пункты перечислений должны либо начинаться с заглавной буквой и заканчиваться точкой, либо начинаться со строчной и закачиваться точкой с запятой; последний пункт пере\-числений всегда заканчивается точкой.
  \item Перед выкладкой финальной версии необходимо просмотреть лог \LaTeX'a, и обратить внимание на сообщения вида \emph{Overfull \textbackslash hbox (1.29312pt too wide) in paragraph}. Обычно, это означает, что текст выползает за поля, и надо подсказать, как правильно слова на слоги, чтобы перенос произошел автоматически. Это делается, например, так: \verb=соломо\-волокуша=.
\end{itemize}



\subsection{Листинги, картинка и прочий \enquote{не текст}}

Различный \enquote{не текст} имеет свойство отображаться не там, где он написан в текстовом в \LaTeX{}, поэтому у него должна быть самодостаточ\-ная понятная подпись \verb=\caption{...}=, уникальная метка \verb=\label{...}=, чтобы на неё можно было бы ссылаться в тексте с помощью \verb=\ref{...}=. Ниже вы сможете увидеть таблицу \ref{time_cmp_obj_func}, на которую мы сослались буквально только что.

\enquote{Не текста} может быть довольно много, чтобы не засорять корневую папку, хорошим решением будет складывать весь \enquote{не текст} в папку \texttt{figures}.
Заклинание \verb=\includegraphics{}= уже знает этот путь и будет искать там файлы без указания папки.
Команда \verb=\input{}= умеет ходит по путям, например \verb=\input{figures/my_awesome_table.tex}=.
Кроме того, листинги кода можно подтягивать из файла с помощью команды \verb=\lstinputlisting{file}=.
%% TODO: Проверить, что на Windows \input{folder/file} работает, если нет, использовать пакет import

%% Вставка кода с помощью listings
\begin{lstlisting}[caption={Название для листинга кода. Достаточно длинное, чтобы люди, которые смотрят картинку сразу после названия статьи (т.~е. все люди), смогли разобраться и понять к чему в статье листинги, картинки и прочий \enquote{не текст}.}, language=Caml, frame=single]
  let x = 5 in x+1
\end{lstlisting}
%% Вставка кода с помощью minted
% \begin{listing}
%   \caption{Название для листинга кода. Достаточно длинное, чтобы люди, которые смотрят картинку сразу после названия статьи (т.~е. все люди), смогли разобраться и понять к чему в статье листинги, картинки и прочий \enquote{не текст}.}
%   \begin{minted}[frame=single]{ocaml}
%     let x = 5 in x+1
%   \end{minted}
% \end{listing}



\subsubsection{Выделения куска листинга с помощью tikz}
Это обывает полезно в текста, а ещё чаще~--- в презентациях. Пример сделан на основе вопроса на \textsc{StackOveflow}\footnote{\url{https://tex.stackexchange.com/questions/284311} (дата доступа:   \DTMdate{2022-12-16}).}.

\begin{figure}
% TODO(Kakadu): Сделать \lstset глобально, чтобы не выписывать все опции листингов каждый раз
\begin{lstlisting}[escapechar=!, basicstyle=\ttfamily, language=c, keepspaces=true]
#include <stdio.h>
#include <math.h>

int main(void)
{
  double c = -1;
  double z = 0;

  (* Это комментарий на русском языке *)
  printf ("For c = %lf:\n", c);
  for (int i=0; i<10; i++ ) {
    printf ( !\tikzmark{a}!"z %d = %lf\n"!\tikzmark{b}!, i, z);
    z = pow(z, 2) + c;
  }
}
\end{lstlisting}

\begin{tikzpicture}[use tikzmark]
\draw[fill=gray,opacity=0.1]
  ([shift={(-3pt,2ex)}]pic cs:a)
    rectangle
  ([shift={(3pt,-0.65ex)}]pic cs:b);
\end{tikzpicture}
\caption{Пример листинга и \textsc{TIKZ} декорации к нему, оформленные в окружении \texttt{figure}. Обратите внимание, что рисунок отображается не там, где он в документе, а может \enquote{плавать}.}
\end{figure}

% !TeX spellcheck = ru_RU
% !TEX root = vkr.tex

\section{Эксперимент (желательно к Новому году)}
Как мы проверяем, что  всё удачно получилось.  К Новому году для промежуточного отчета желательно хотя бы описать как он будет прово\-диться и на чем.

\subsection{Условия эксперимента}
Железо (если актуально);  версии ОС, компиляторов и параметры командной строки; почему мы выбрали именно эти тесты; входные дан\-ные, на которых проверяем наш подход, и почему мы выбрали именно их.

\subsection{Исследовательские вопросы }
По-английски называется \emph{research questions}, в тексте можно ссылаться на них как RQ1, RQ2, и т.~д.
Необходимо сформулировать, чего мы хотели бы добиться работой (2 пункта будет хорошо):

\begin{itemize}
\item Хотим алгоритм, который лучше вот таких-то остальных.
\item Если в подходе можно включать/выключать составляющие, то насколько существенно каждая составляющая влияет на улучшения.
\item Если у нас строится приближение каких-то штук, то на сколько точными будут эти приближения.
\item и т.п.
\end{itemize}

Иногда в работах это называют гипотезами, которые потом проверяют. Далее в тексте можно ссылаться на research questions как \textsc{RQ}, это обще\-при\-нятое сокращение.

\subsection{Метрики}

Как мы сравниваем, что результаты двух подходов лучше или хуже:
\begin{itemize}
\item Производительность.
\item Строчки кода.
\item Как часто алгоритм \enquote{угадывает} правильную класси\-фикацию входа.
\end{itemize}

\noindent Иногда метрики вырожденные (да/нет), это не очень хорошо, но если в области исследований так принято, то ладно.

\subsection{Результаты}
Результаты понятно что такое. Тут всякие таблицы и графики, как в таблице \ref{time_cmp_obj_func}. Обратите внимание, как цифры выровнены по правому краю, названия по центру, а разделители $\times$ и $\pm$ друг под другом.

Скорее всего Ваши измерения будут удовлетворять нормальному распределению, в идеале это надо проверять с помощью критерия Кол\-могорова и т.п.
Если критерий этого не подтверждает, то у Вас что-то сильно не так с измерениями, надо проверять кэши процессора, отключать Интернет во время измерений, подкручивать среду исполне\-ния (англ. runtime), что\-бы сборка мусора не вмешивалась и т.п.
Если критерий удовлетворён, то необходимо либо указать мат. ожидание и доверительный/предсказы\-вающий интервал, либо написать, что все измерения проводились с погрешностью, например, в 5\%.
Замечание: если у вас получится улуч\-шение производительности в пределах погреш\-ности, то это обязательно вызовет вопросы.

В этом разделе надо также коснуться Research Questions.

\subsubsection{RQ1} Пояснения
\subsubsection{RQ2} Пояснения

\begin{table}
\def\arraystretch{1.1}  % Растяжение строк в таблицах
\setlength\tabcolsep{0.2em}
\centering
% \resizebox{\linewidth}{!}{%
    \caption{Производительность какого-то алгоритма при различных разрешениях картинок  (меньше~--- лучше), в мс.,  CI=0.95. За пример таблички кидаем чепчики в честь Я.~Кириленко}
    \begin{tabular}[C]{
    S[table-format=4.4,output-decimal-marker=\times]
    *4{S
          [table-figures-uncertainty=2, separate-uncertainty=true, table-align-uncertainty=true,
          table-figures-integer=3, table-figures-decimal=2, round-precision=2,
          table-number-alignment=center]
          }
    }
    \toprule
        \multicolumn{1}{r}{Resolution} & \multicolumn{1}{r}{\textsc{TENG}} & \multicolumn{1}{r}{\textsc{LAPM}} &
        \multicolumn{1}{r}{\textsc{VOLL4}} \\ \midrule
        1920.1080 & 406.23 \pm 0.94 & 134.06 \pm 0.35 & 207.45 \pm 0.42  \\ \midrule
        1024.768  & 145.0 \pm 0.47  & 39.68 \pm 0.1   &  52.79  \pm 0.1 \\ \midrule
        464.848   & 70.57 \pm 0.2   & 19.86 \pm 0.01     & 32.75  \pm 0.04 \\ \midrule
        640.480   & 51.10 \pm 0.2   & 14.70 \pm 0.1 & 24  \pm 0.04 \\ \midrule
        160.120   & 2.4 \pm 0.02    & 0.67 \pm 0.01      & 0.92  \pm 0.01 \\
        \bottomrule
    \end{tabular}%
%}
    \label{time_cmp_obj_func}
\end{table}

\clearpage
% !TeX spellcheck = ru_RU
% !TEX root = vkr.tex

\newcolumntype{C}{ >{\centering\arraybackslash} m{4cm} }
\newcommand\myvert[1]{\rotatebox[origin=c]{90}{#1}}
\newcommand\myvertcell[1]{\multirowcell{5}{\myvert{#1}}}
\newcommand\myvertcelll[1]{\multirowcell{4}{\myvert{#1}}}
\newcommand\myvertcellN[2]{\multirowcell{#1}{\myvert{#2}}}

\afterpage{%
	\clearpage% Flush earlier floats (otherwise order might not be correct)
	\thispagestyle{empty}% empty page style (?)
	\begin{landscape}% Landscape page
		\centering % Center table

		\begin{tabular}{|c|c|c|c|c|c|c|c|c|c|c|c|c|c|c|c|c|c|}\hline
			%& \multicolumn{17}{c|}{} \\ \hline
			\multirowcell{2}{Код модуля \\в составе \\ дисциплины,\\практики и т.п. }
			&\myvertcellN{2}{Трудоёмкость\quad}
			& \multicolumn{10}{c|}{\tiny{Контактная работа обучающихся с преподавателем}}
			& \multicolumn{5}{c|}{\tiny{Самостоятельная работа}}
			& \myvertcellN{2}{\tiny Объем активных и интерактивных\quad}
			\\ \cline{3-17}

			&& \myvertcellN{2}{лекции\quad}
			&\myvertcellN{2}{семинары\quad}
			&\myvertcellN{2}{консультации\quad}
			&\myvertcellN{2}{\small практические  занятия\quad}
			&\myvertcellN{2}{\small лабораторные работы\quad}
			&\myvertcellN{2}{\small контрольные работы\quad}
			&\myvertcellN{2}{\small коллоквиумы\quad}
			&\myvertcellN{2}{\small текущий контроль\quad}
			&\myvertcellN{2}{\small промежуточная аттестация\quad}
			&\myvertcellN{2}{\small итоговая аттестация\quad}

			&\myvertcellN{2}{\tiny под руководством    преподавателя\quad}
			&\myvertcellN{2}{\tiny в присутствии     преподавателя\quad}
			&\myvertcellN{2}{\tiny с использованием    методических\quad}
			&\myvertcellN{2}{\small текущий контроль\quad}
			&\myvertcellN{2}{\makecell{\small промежуточная \\ аттестация}}
			&     \\
			&& &&&&&&&&& &&&&&&\\
			&& &&&&&&&&& &&&&&&\\
			&& &&&&&&&&& &&&&&&\\
			&&&&&&&&&&& &&&&&&\\
			&&&&&&&&&&& &&&&&&\\
			&&&&&&&&&&& &&&&&&\\ \hline
			Семестр 3 & 2 &30  &&&&&&&&2   & &&&18 &&20 &10\\ \hline
			&   &2--42&&&&&&&&2--25& &&&1--1&&1--1&\\ \hline
			Итого     & 2 &30  &&&&&&&&2   & &&&18 &&20 &10\\ \hline
		\end{tabular}

		\captionof{table}{Если таблица очень большая, то можно её изобразить на отдельной портретной странице. Не забудьте подробное описание, чтобы содержимое таблицы можно было понять не читая весь текст.}
	\end{landscape}
	\clearpage% Flush page
}



\subsection{Обсуждение результатов}

Чуть более неформальное обсуждение, то, что сделано. Например, почему метод работает лучше остальных? Или, что делать со случаями, когда метод классифицирует вход некорректно.

% !TeX spellcheck = ru_RU
% !TEX root = vkr.tex

\section{Применение}

Применение того, что сделано на практике. Например, оно может не работать в лоб, потому что всё изложено чуть более сжато и теоретично, надо рассказать тонкости и правильный метод применения результатов. Если результаты применяются без до\-полнительных телодвижений, то про это можно не писать.

% !TeX spellcheck = ru_RU
% !TEX root = vkr.tex

\section{Угрозы нарушения корректности (опциональный)}

Если основная заслуга метода, это то, что он дает лучшие цифры, то стоит сказать, где мы могли облажаться, когда
\begin{enumerate}
  \item проводили численные замеры;
  \item выбирали тестовый набор (см. \emph{confirmation bias}).
\end{enumerate}

% !TeX spellcheck = ru_RU
% !TEX root = vkr.tex

\section{Реализация}
Очень важный раздел для будущих программных инженеров, т.е. почти для всех. Важно иметь всегда, в том числе для промежуточных отчетов по учебным практикам или ВКР.

В процессе работы можно сделать огромное количество косяков, неполный список которых ниже.

\begin{enumerate}
  \item Реализация должна быть. На публично доступную реализацию обязательная ссылка. Если код под \textsc{NDA}, то об этом,
  во-первых, должно быть сказано явно,
  во-вторых, на защиту должны выно\-ситься другие результаты (например, архитектура), чтобы комис\-сия имела возможность оценить хоть что-то,
  и, в третьих, должна быть справка от работодателя, что вы правда что-то сделали.
        \begin{itemize}
          \item  Рецензент обязан оценить код (о возможности должен побеспо\-коиться обучающийся).
        \end{itemize}
  \item Код реализации должен быть написан защищающимся целиком.
        \begin{itemize}
          \item  Если проект групповой, то нужно явно выделить какие части были модифицированы защищающимся. Например, в преды\-дущих разделах на картинке архитектуры нужно выделить цветом то, что вы модифицировали.
          \item Нельзя пускать в негрупповой проект коммиты от других людей, или людей не похожих на Вас. Например, в 2022 году защищающийся-парень делал коммиты от сценического псев\-донима, который намекает на женский \enquote{гендер}. (Нет, это не шутка.) На тот момент в российской культуре это выглядело странно.
          \item Возможна ситуация, что вы используете конкретный ник в интернете уже лет пять, и желаете писать ВКР под этим ником на \GitHub{}. В принципе, это допустимо (не только лишь я так считаю), но если Вы встретите преподавателя, который считает наоборот, то Вам придется грамотно отмазы\-ваться. В Вашу пользу могут сыграть те факты, что к нику на гитхабе у Вас приписаны настоящие имя и фамилия; что в репозитории у вас видна домашка за первый курс; и что Ваш преподаватель практики сможет подтвердить, что Вы уже несколько лет используете это ник; и т.п.
        \end{itemize}
  \item Если вы получаете диплом о присвоении звания программного инженера, код должен соответствовать.
        \begin{enumerate}
          \item Не стоит выкладывать код одним коммитом.
          \item Не стоит выкладывать код аккурат перед защитой.
          \item Лучше хоть какие-то тесты, чем совсем без них. В идеале нужно предъявлять процент покрытия кода тестами.
          \item Лучше  сделать \textsc{CI}, а также \textsc{CD}, если оно уместно в Вашем проекте.
          \item Не стоит демонстрировать на защите, что Вам даже не пришло в голову напустить на код линтеры и т.п.
        \end{enumerate}
  \item Если ваша реализация по сути является прохождением стандартного туториала,
  например, по отделению картинок кружек от котиков с помощью машинного обучения, то необходимо срочно сообщить об этом куратору на мат-мехе,
  иначе Государственная Экзаменацион\-ная Комиссия \enquote{порвёт Вас как Тузик грелку}, поставит \enquote{единицу},
  а все остальные Ваши сокурсники получат оценку выше. (Это не шутка, а реальная история 2020 года.)
\end{enumerate}

\noindent Если Вам предстоит защищать учебную практику, а эти рекомендации видятся как более подходящие для защиты ВКР, то ... отмаза не засчиты\-вается, сразу учитесь делать нормально.

% !TeX spellcheck = ru_RU
% !TEX root = vkr.tex

\section*{Заключение}
\textbf{Обязательно для промежуточного, полугодового, годового и  любых других отчётов.} Кратко, что было сделано.

\textbf{Для практик/ВКР.} Также важно сделать список результатов, который будет один к одному соответствовать задачам из раздела~\ref{sec:task}.

\begin{itemize}
\item Результат к задаче №1.
\item Результат к задаче №2.
\item и т.д.
\end{itemize}
\noindent Для промежуточных отчетов сюда важно записать какие задачи уже были сделаны за осенний семестр, а какие только планируется сделать.

Также сюда можно написать планы развития работы в будущем, или, если их много, выделить под это отдельную предпоследнюю главу.



\setmonofont{CMU Typewriter Text}
\bibliographystyle{ugost2008ls}
\bibliography{vkr}
\end{document}
